\documentclass[aspectratio=169,11pt]{beamer}

\usetheme{Madrid}
\usecolortheme{default}
\usefonttheme{professionalfonts}
\setbeamertemplate{navigation symbols}{}
\setbeamertemplate{footline}[frame number]

\usepackage{amsmath, amssymb, booktabs, graphicx}

\title{Comparative and Global Labor-Market Institutions}
\subtitle{Labor Markets and Political/Cultural Institutions -- Week 9}
\author{Teaching deck draft}
\date{}

\begin{document}

\begin{frame}
  \titlepage
\end{frame}

\begin{frame}{Central question}
\begin{itemize}
    \item How do labor-market institutions differ across development settings, welfare states, migration regimes, and global production systems?
    \item Which lessons travel across contexts, and which estimates are setting-specific?
    \item The labor margins are wages, employment, unemployment, formality, protection, integration, and working conditions.
\end{itemize}
\end{frame}

\begin{frame}{Distinction from Week 8}
\begin{itemize}
    \item Week 8: how reforms work once implementation, knowledge, and adjustment are uneven.
    \item Week 9: how entire institutional regimes differ before any one reform is studied.
    \item Reform analysis asks what changed.
    \item Comparative analysis asks what setting, bundle, and governance level condition the effect.
\end{itemize}
\end{frame}

\begin{frame}{Comparative taxonomy}
\small
\begin{tabular}{p{0.25\linewidth} p{0.28\linewidth} p{0.34\linewidth}}
\toprule
Dimension & Core question & Labor margin \\
\midrule
Formal rules & What does law or policy state? & rights, costs, procedures \\
Coverage & Who is legally reached? & covered jobs, firms, worker groups \\
Membership & Who belongs to an organization or system? & union, insurance, migration status \\
Enforcement & Is the rule credible? & compliance, formality, claims \\
Complements & What other institutions move with it? & wage-setting, security, adjustment \\
Governance level & Who regulates? & state, firm, buyer, transnational body \\
\bottomrule
\end{tabular}
\end{frame}

\begin{frame}{Comparative regime framework}
\small
\[
Y^{labor}_{c,s}
= g(\text{rules}_{c}, \text{coverage}_{c,s}, \text{enforcement}_{c,s},
\text{complements}_{c}, \text{global exposure}_{c,s})
\]
\begin{itemize}
    \item Countries, sectors, firms, and worker groups can face different effective institutions.
    \item Similar legal text can produce different labor outcomes when coverage or enforcement differs.
    \item Mechanism portability is easier than estimate portability.
\end{itemize}
\end{frame}

\begin{frame}{Regime domains}
\small
\begin{tabular}{p{0.30\linewidth} p{0.30\linewidth} p{0.28\linewidth}}
\toprule
Domain & Institution & Observed margin \\
\midrule
Welfare states & UI, EPL, active labor policy & unemployment, job-finding \\
Bargaining systems & coverage, coordination, extension & wages, rents, inequality \\
Labor regulation indices & employment, collective, social law & rules on paper, protection \\
Informality regimes & registration, tax, inspection & formality, earnings, compliance \\
Migration regimes & entry, rights, recognition & integration, matching, firm response \\
Global production & codes, audits, buyers & conditions, compliance, turnover \\
\bottomrule
\end{tabular}
\end{frame}

\begin{frame}{Welfare states and complementarities}
\begin{itemize}
    \item Blanchard--Wolfers: shocks and institutions interact in unemployment dynamics.
    \item Observed margin: aggregate unemployment across countries and time.
    \item Core comparative point: institutions are bundles, not one-dimensional rankings.
    \item Welfare interpretation needs security, job-finding, wage-setting, and incidence.
\end{itemize}
\end{frame}

\begin{frame}{Bargaining systems}
\begin{itemize}
    \item Union membership, bargaining coverage, agreement extension, and coordination are distinct.
    \item Membership measures organization; coverage measures exposure to wage-setting rules.
    \item Industrial relations systems affect wages, compression, firm rents, and adjustment to shocks.
    \item A country comparison must name the actual exposure margin.
\end{itemize}
\end{frame}

\begin{frame}{Cross-country regulation indices}
\begin{itemize}
    \item Botero--Djankov--La Porta--Lopez-de-Silanes--Shleifer: cross-country labor-regulation coding.
    \item Identifying variation: legal and institutional differences across countries.
    \item Observed object: rules-on-paper indices and institutional correlates.
    \item Main limit: legal text is not implementation, enforcement, or worker coverage.
\end{itemize}
\end{frame}

\begin{frame}{Informality and development settings}
\begin{itemize}
    \item Informality is an institutional regime, not a missing-data category.
    \item Ulyssea: informality involves firm registration, payroll, inspection, worker protection, and productivity.
    \item Observed margins: formality, earnings, firm size, compliance, productivity.
    \item Comparative caution: the same legal protection can cover very different shares of work.
\end{itemize}
\end{frame}

\begin{frame}{Migration regimes and integration}
\begin{itemize}
    \item Migration rules govern entry, work authorization, rights, credential recognition, and benefit portability.
    \item Beerli--Peri--Ruffner--Siegenthaler: Swiss opening to cross-border workers.
    \item Identifying variation: differential exposure to a migration-regime change.
    \item Observed margins: firm performance, employment, worker outcomes, matching.
\end{itemize}
\end{frame}

\begin{frame}{Global supply chains and private regulation}
\begin{itemize}
    \item Labor governance can be transnational: buyers, codes, audits, certification, and supplier management.
    \item Locke: private power can complement public regulation but has limits.
    \item Distelhorst--Hainmueller--Locke: Nike supply-chain evidence on management and labor standards.
    \item Observed margins: compliance, working conditions, social performance, turnover.
\end{itemize}
\end{frame}

\begin{frame}{Making country evidence travel}
\small
\begin{tabular}{p{0.29\linewidth} p{0.58\linewidth}}
\toprule
Move & What to do \\
\midrule
Mechanism first & Lead with wage-setting, informality, integration, bargaining, or conditions \\
Taxonomy & Place the country in a welfare, bargaining, informality, migration, or supply-chain regime \\
Leverage & State the reform, border, discontinuity, administrative data, or firm linkage \\
Observed margin & Name the labor outcome actually measured \\
Transportability & Say what should travel and what depends on local complements \\
\bottomrule
\end{tabular}
\end{frame}

\begin{frame}{Critiques of comparative evidence}
\begin{itemize}
    \item The setting may be too idiosyncratic.
    \item Institutional bundles may confound the focal mechanism.
    \item Legal categories and outcome measures may not be comparable across countries.
    \item Equilibrium feedback can differ across labor, product, and political markets.
    \item Rules on paper can be mistaken for effective implementation.
\end{itemize}
\end{frame}

\begin{frame}{Empirical design menu}
\small
\begin{tabular}{p{0.33\linewidth} p{0.28\linewidth} p{0.28\linewidth}}
\toprule
Design & Identifying variation & Observed margin \\
\midrule
Cross-country index & legal coding across countries & regulation, protection \\
Within-country reform & timing, threshold, rollout & wages, formality, unemployment \\
Border or region exposure & migration or market access & employment, matching, firms \\
Administrative microdata & regime differences in linked records & earnings, mobility, benefits \\
Supply-chain data & audit, buyer, certification exposure & compliance, conditions \\
Historical classification & inherited institutional regimes & modern labor outcomes \\
\bottomrule
\end{tabular}
\end{frame}

\begin{frame}{What travels}
\begin{itemize}
    \item Mechanisms: enforcement credibility, coverage gaps, informality margins, migration-induced matching, buyer incentives.
    \item Design logic: what variation identifies and which labor margin is observed.
    \item Conceptual distinctions: rules versus implementation, coverage versus membership, public versus private governance.
    \item Comparative taxonomies: regime, sector, border, firm network, or supply chain.
\end{itemize}
\end{frame}

\begin{frame}{What does not travel mechanically}
\begin{itemize}
    \item Point estimates when capacity, bargaining, informality, market structure, or politics differ.
    \item Welfare incidence when costs shift across workers, firms, prices, profits, and entry.
    \item Private regulation effects when buyer commitment or worker voice differs.
    \item Migration effects when credential rules, firm technology, or local labor supply differs.
\end{itemize}
\end{frame}

\begin{frame}{Week 9 lab}
\begin{itemize}
    \item Reproduce: construct a synthetic labor-regulation index inspired by Botero et al.
    \item Diagnose: classify implementation gaps and transportability risks.
    \item Transfer: map designs to migration, informality, welfare-state, and supply-chain settings.
    \item Output: index tables, transportability diagnostic, transfer design classification.
\end{itemize}
\end{frame}

\begin{frame}{Bridge to Week 10}
\begin{itemize}
    \item Week 9 gives the comparative vocabulary for synthesis.
    \item Week 10 turns the course into student research designs.
    \item A strong design names the labor outcome, institutional mechanism, comparison class, identifying variation, and portability claim.
    \item Institutional detail should sharpen the labor question, not replace it.
\end{itemize}
\end{frame}

\begin{frame}{Takeaways}
\begin{itemize}
    \item Comparative labor evidence changes how we interpret institutional effects.
    \item Similar formal rules can operate differently because coverage, enforcement, complements, and global exposure differ.
    \item Migration and supply chains extend institutional comparison beyond closed national labor markets.
    \item Country-specific evidence matters when it teaches a portable mechanism with honest limits.
\end{itemize}
\end{frame}

\end{document}
